\documentclass[12pt]{article}
% Préambule commun à tous les documents extraits de « La Revue CAPES »
\usepackage[T1]{fontenc}
\usepackage[utf8]{inputenc}
\usepackage{lmodern}
\usepackage{amsmath,amssymb,amsthm,amsfonts,mathrsfs}
\usepackage{setspace}
\usepackage{listings}
\usepackage{pgfplots}
\pgfplotsset{compat=1.18}
\usepackage{longtable}
\usepackage{adjustbox}
\usepackage{lettrine}
\usepackage{tkz-tab}
\usepackage{changepage}
\usepackage{pdflscape}
\usepackage{graphicx}
\usepackage{tikz}
\usepackage{xcolor}
\usepackage[a4paper,top=2cm,bottom=2.5cm,left=2.5cm,right=2cm]{geometry}
\usepackage{fancyhdr}
\usepackage{draftwatermark}
\SetWatermarkText{La RevueCapes}
\SetWatermarkLightness{0.9}
\SetWatermarkScale{2}
\usepackage{hyperref}
\hypersetup{colorlinks=true,linkcolor=blue!50!black,urlcolor=blue!50!black}

\definecolor{revue}{RGB}{20,60,120}

\pagestyle{fancy}
\fancyhf{}
\renewcommand{\headrulewidth}{0.4pt}
\setlength{\headheight}{15pt}

% \entete{Matière}{Titre}{Type de document}
\newcommand{\entete}[3]{%
  \fancyhead[L]{\small\textsc{La Revue CAPES} -- #1}%
  \fancyhead[R]{\small #2}%
  \fancyfoot[C]{\thepage}%
  \thispagestyle{plain}%
  \begin{center}
    {\color{revue}\rule{\textwidth}{1.2pt}}\\[4pt]
    {\small\textsc{Concours directs CAP-CEG / CAPES -- Burkina Faso}}\\[6pt]
    {\LARGE\bfseries\color{revue} #2}\\[6pt]
    {\large\itshape #1 -- #3}\\[2pt]
    {\color{revue}\rule{\textwidth}{1.2pt}}
  \end{center}
  \vspace{0.5cm}%
}

% Encadré pour les remarques ajoutées dans les corrigés
\newcommand{\remarque}[1]{\par\medskip\noindent\fcolorbox{revue}{revue!6}{\parbox{\dimexpr\linewidth-2\fboxsep-2\fboxrule}{\textbf{\color{revue}Remarque.} #1}}\par\medskip}


\begin{document}
\entete{Mathématiques}{Sujet type n°4}{Énoncé}
\begin{spacing}{1.5}

 
 \textbf{Exercice 1}
 
  On considère l'ensemble $U$ des nombres complexes de module égal à 1. Soit $a$, un nombre
 complexe tel que $\mid a\mid= 1$.
 
 1. Montrer que l'application $f_a$ donnée par :
 
 $f_a(z) = \frac{z+a}{1+\overline{a}z} $
 
 est bien définie pour tout élément $z$ de $U$.
 
 2. Montrer que, si $z\in U$, alors $\overline{z} = \frac{1}{z}$.
 
 3. En déduire que si $z\in U$, alors $f_a(z)\in U$.
 
 4. Réciproquement, montrer que tout élément $t$ de $U$ est l'image par $f_a$ d'un unique élément
 $z$ de $U$ que l'on déterminera.
 
 5. Déduire de ce qui précède que $f_a$ est une bijection de $U$ sur $U$ et préciser sa bijection
 réciproque.
 
 6. Donner l'ensemble des points $z$ dont l'image par $f_a$ appartient à l'ensemble $\lbrace {-1,1,i,-i} \rbrace$ dans chacun des cas suivants :
 
 (a) $a = 2$;
 
 (b) $a = 2i$;
 
 (c) $a = 1+i$.
 
 \textbf{Exercice 2 }
 
  Soit $a$ un paramètre réel. On considère la forme quadratique sur $\mathbb{R}^3$ définie par
 
 $q_a((x_1,x_2,x_3)) = x^{2}_1 + (1 + a)x^{2}_2 + (1 + a+a^2)x^{2}_3 +2x_1x_2-2ax_2x_3$.
 
 1. Donner la matrice de $q_a$ dans la base canonique de $\mathbb{R}^3$.
 
 2. Déterminer la $f.b.s$ associée à $q_a$.
 
 3. Donner la décomposition de Gauss de $q_a$.
 
 4. Donner le noyau, le rang et la signature de $q_a$ en fonction suivant les valeurs de $a$.
 
 5. Déterminer une base $q_2-orthogonale$ de $\mathbb{R}^3$ notée $B$.
 
 6. Pour quelles valeurs de $a$, $q_a$ est définie positive.
 
 \medskip\textbf{Exercice 3}
 
  Soient
 
 $I = \int_{0}^{\frac{\pi}{2}}\frac{\cos^2(x)}{\cos(x) + \sin(x)}dx$ et $J = \int_{0}^{\frac{\pi}{2}}\frac{\sin^2(x)}{\cos(x) + \sin(x)}dx$ 

 1. Sans calculer $I$ et $J$ montrer que $I = J$
 
 2. Vérifier que $\forall {x\in \mathbb{R}}, \cos(x)+\sin(x) = 2\cos(x-\frac{\pi}{4})$.
 
 3. En déduire que $I+J = \frac{\sqrt{2}}{2}\int_{\frac{-\pi}{4}}^{\frac{\pi}{4}}\frac{dx}{\cos(x)}$.
 
 4) Calculer $I + J$. En déduire $I$ et $J$.
 
 \textbf{Exercice 4 }
 
  1. Trouver les valeurs des constantes $a$ et $b$ telles que :

$\frac{x^2}{\sqrt{x^2-1}}=\frac{a}{\sqrt{x^2-1}} + \frac{\sqrt{x^2-1}}{b}$

 2. En déduire l'aire de la surface délimitée par la courbe d'équation $y = \frac{x^3}{\sqrt{x^2-1}}$, l'axe des $x$ et les droites d'équations $x = \sqrt{2}$ et $x = 2$.
 
 
 \quad
 \quad

\end{spacing}
\end{document}
